\documentclass[10pt]{article}

\usepackage[margin=1in]{geometry}
\usepackage{booktabs}
\usepackage{hyperref}
\usepackage{enumitem}
\usepackage{xcolor}

\hypersetup{
  colorlinks=true,
  linkcolor=blue,
  urlcolor=blue,
  citecolor=blue
}

\title{RV-MalTrace: Hardware-Assisted Behavioral Tracing for RISC-V Linux Workloads}
\author{Anonymous Authors}
\date{NDSS 2026 Draft Artifact Skeleton}

\newcommand{\artifactroot}{\texttt{results/evaluation/genesys2-cva6/current/}}
\newcommand{\todo}[1]{\textcolor{red}{TODO: #1}}

\begin{document}
\maketitle

\begin{abstract}
RV-MalTrace is a CVA6 and Genesys2 trace path for controlled RISC-V Linux behavior studies.
The current artifact package supports bounded hardware trace claims for safe synthetic workloads, field-level semantic provenance, and reproducible local checkers rooted at \artifactroot.
The paper draft explicitly separates hardware-derived fields from exact-ELF, runtime-map, and oracle-derived fields.
It does not claim real-malware validation, production streaming/DMA throughput, board-native DWARF source lines from sidecar data, or full hardware pointer-string recovery unless an external intake record accepts those artifacts.
It also does not claim JTAG RAM boot, SD-card kernel update, or board cycle-overhead results while those host/board gates remain blocked.
\end{abstract}

\section{Scope And Claims}
The submission scope is a hardware-assisted behavioral trace system for RISC-V Linux workloads.
All paper-facing evidence is tied to \artifactroot{} and checked by recursive path and SHA-256 validation.
QEMU, strace, and host-side control logs are validation oracles only; they are not reported as hardware recovery output.

\begin{table}[h]
\centering
\caption{Claim boundaries for the NDSS draft.}
\begin{tabular}{p{0.27\linewidth}p{0.24\linewidth}p{0.39\linewidth}}
\toprule
Area & Status & Boundary \\
\midrule
Canonical evidence root & PASS & \artifactroot{} is the only current evidence root. \\
CVA6 trace correctness & PASS local and host xsim & Strict syscall pairing and marker-window ordering are checked by RTL tests. \\
Semantic provenance & PASS & Fields are tagged as hardware, exact ELF, runtime map, or oracle. \\
P0 Genesys2 BRAM evidence & PASS & Controlled safe synthetic marker windows only. \\
Local code attribution & PASS local fixtures & Exact ELF, PIE/ASLR, dynamic object, fork/exec, stripped-ELF, and sidecar-boundary fixtures. \\
JTAG RAM boot & BLOCKED & Current target exposes ILA but no Vivado memory-control object. \\
Production streaming/DMA & BLOCKED & No throughput claim without accepted external artifacts. \\
Real malware & NOT CLAIMED & Safe malware-like workloads are not real malware validation. \\
\bottomrule
\end{tabular}
\end{table}

\section{Trace Architecture}
The CVA6 adapter emits event records for retire, trap, syscall entry, syscall return, privilege transition, branch, jump, context, pointer-byte snapshots, and markers.
The paper configuration keeps relaxed SRET qualification disabled.
BRAM marker windows clear on begin markers, freeze on end markers, and record the begin marker as sequence zero even when the previous window was frozen.

\section{Semantic Reconstruction}
The analysis pipeline combines hardware trace records with exact board ELF metadata and runtime process maps.
Static analysis supports exact ELF SHA-256 checks, PIE and ASLR load bias, dynamic-object attribution, fork/exec ownership, and stripped-ELF boundaries.
Oracle traces from QEMU or strace are used to validate expected behavior and must not be substituted for hardware-derived semantics.
Sidecar-derived source maps are never counted as board-native DWARF unless an accepted external summary binds them to exact board artifacts.

\section{Evaluation Plan}
\begin{table}[h]
\centering
\caption{Current evaluation table.}
\begin{tabular}{p{0.16\linewidth}p{0.34\linewidth}p{0.18\linewidth}p{0.22\linewidth}}
\toprule
RQ & Measurement & Status & Checker \\
\midrule
RQ1 & Trace ordering, trap/retire separation, syscall pairs & PASS & \texttt{check\_fuzz\_trace\_plan.py} \\
RQ2 & Field provenance and semantic reconstruction & PASS & \texttt{check\_genesys2\_semantic\_provenance.py} \\
RQ3 & Exact ELF, PIE/ASLR, runtime map, dynamic-object, fork/exec, and stripped-ELF attribution & PASS local fixtures & \texttt{check\_genesys2\_local\_code\_analysis\_fixtures.py} \\
RQ4 & Production streaming/DMA transport & BLOCKED & external intake \\
RQ5 & JTAG RAM boot or SD-card-free kernel update & BLOCKED & \texttt{check\_genesys2\_jtag\_ram\_boot\_probe.py} \\
\bottomrule
\end{tabular}
\end{table}

\section{Artifact Instructions}
Quick local reproduction is:
\begin{verbatim}
uv run rvmt repro:quick
\end{verbatim}
Docker-local reproduction is:
\begin{verbatim}
uv run rvmt ndss:docker-full
\end{verbatim}
Host-only Vivado and board work is documented in \texttt{docs/07-evaluation-evidence/ndss\_host\_runbook.md}.
Host LaTeX compilation for this draft is:
\begin{verbatim}
latexmk -pdf -interaction=nonstopmode -halt-on-error \
  -outdir=build/latex/ndss2026 docs/08-publication/ndss2026/paper.tex
\end{verbatim}

\section{Limitations}
The current artifact package is intentionally conservative.
Production streaming/DMA throughput and board cycle-overhead remain blocked until their summaries are accepted with real, hashed artifacts.
The current JTAG RAM-boot probe is read-only and blocked because the visible hardware objects do not expose a RAM-write or hart-control path.
No PASS status is assigned to Vivado, board, LaTeX, or malware experiments that were not actually run.

\section{Ethics And Safety}
The repository contains controlled safe synthetic and malware-like workloads for measurement.
It does not provide real malware execution evidence in the current NDSS claim set.
Any future real-malware handling must use containment, lineage, and release-policy gates before being cited.

\section{Conclusion}
The current draft positions RV-MalTrace as a reproducible, claim-bounded RISC-V hardware tracing artifact.
The next paper task is to replace this skeleton with the official NDSS template and fill in narrative, figures, and citations without widening claims beyond the accepted evidence root.

\end{document}
